%% =======================================================
%% Actividad 31: "Protección de datos personales vs. datos sensibles (resumen)"
%% =======================================================

%% ---- CONFIGURACIÓN ------------------------

\documentclass[11pt,letterpaper]{article}

\usepackage[indice]{tecnm}
\usepackage{float}
\usepackage{booktabs}
\usepackage[T1]{fontenc}
\usepackage{xcolor}
\usepackage[most]{tcolorbox} % Carga tcolorbox con soporte para listings y saltos de página
\tcbuselibrary{listings, skins, breakable}

% Paleta de colores (estilo GitHub Light / VS Code)
\definecolor{codebg}{HTML}{F8F9FA}       % Fondo del código
\definecolor{codehead}{HTML}{E9ECEF}     % Fondo de la barra de título
\definecolor{codeborder}{HTML}{CED4DA}   % Borde de la tarjeta
\definecolor{codetitle}{HTML}{212529}    % Texto del título
\definecolor{codekw}{HTML}{0550AE}       % Palabras clave (import, def, for...)
\definecolor{codefunc}{HTML}{8250DF}     % Funciones y métodos (describe, groupby...)
\definecolor{codestr}{HTML}{0A3069}      % Strings
\definecolor{codecomm}{HTML}{6E7781}     % Comentarios
\definecolor{codenum}{HTML}{8C959F}      % Números de línea
\definecolor{codeid}{HTML}{24292F}       % Texto general

% Configuración base de listings
\lstdefinestyle{pythonestilo}{
  language=Python,
  basicstyle=\ttfamily\footnotesize\color{codeid},
  keywordstyle=\color{codekw}\bfseries,
  keywordstyle=[2]\color{codefunc}, % Estilo secundario para métodos de datos
  stringstyle=\color{teal!80!black},
  commentstyle=\color{codecomm}\itshape,
  morekeywords={as, True, False, None, with, yield, lambda},
  morekeywords=[2]{read_csv, head, tail, info, describe, groupby, mean, value_counts, to_csv, print, len, range},
  numbers=left,
  numberstyle=\tiny\ttfamily\color{codenum},
  stepnumber=1,
  numbersep=10pt,
  showstringspaces=false,
  morestring=[b]{"""},
  tabsize=4,
  breaklines=true,
  breakatwhitespace=true,
  columns=flexible,
  keepspaces=true,
  literate=
    {á}{{\'a}}1 {é}{{\'e}}1 {í}{{\'i}}1 {ó}{{\'o}}1 {ú}{{\'u}}1
    {Á}{{\'A}}1 {É}{{\'E}}1 {Í}{{\'I}}1 {Ó}{{\'O}}1 {Ú}{{\'U}}1
    {ü}{{\char252}}1 {«}{{\guillemotleft}}1 {»}{{\guillemotright}}1 {·}{{\textperiodcentered}}1 {ñ}{{\~n}}1 {→}{{$\rightarrow$}}1 {Ñ}{{\~N}}1 {¿}{{?`}}1 {¡}{{!`}}1
}

% Definición del entorno personalizado \begin{codigo}{Título}{etiqueta}
% \newtcblisting[auto counter, number within=section]{codigo}[2][]{
\newtcblisting[auto counter]{codigo}[2][]{
  listing only,
  breakable,                    % Permite que el bloque se divida si pasa de página
  enhanced,
  colback=codebg,
  colframe=codeborder,
  colbacktitle=codehead,
  coltitle=codetitle,
  fonttitle=\bfseries\small\sffamily,
  boxrule=0.6pt,
  arc=4pt,                      % Bordes redondeados
  left=20pt,                    % Espacio izquierdo para los números de línea
  right=8pt,
  top=6pt,
  bottom=6pt,
  toptitle=4pt,
  bottomtitle=4pt,
  before skip=12pt,
  after skip=12pt,
  listing options={style=pythonestilo},
  title={Código \thetcbcounter: #2},
  #1
}

% Paleta para salida de consola (estilo terminal limpia)
\definecolor{termbg}{HTML}{212529}      % Fondo oscuro tipo terminal
\definecolor{termhead}{HTML}{343A40}    % Barra superior
\definecolor{termtext}{HTML}{F8F9FA}    % Texto claro de alta legibilidad
\definecolor{termborder}{HTML}{495057}  % Borde

% Entorno \begin{consola}[label=out:iris]{Título}
\newtcblisting[auto counter]{consola}[2][]{
  listing only,
  breakable,
  colback=termbg,
  colframe=termborder,
  colbacktitle=termhead,
  coltitle=termtext,
  colupper=termtext,
  fonttitle=\bfseries\small\sffamily,
  boxrule=0.6pt,
  arc=4pt,
  left=8pt,
  right=8pt,
  top=6pt,
  bottom=6pt,
  toptitle=4pt,
  bottomtitle=4pt,
  before skip=12pt,
  after skip=12pt,
  listing options={
    basicstyle=\ttfamily\scriptsize\color{termtext}, % \scriptsize evita que las tablas de Pandas se partan de línea
    breaklines=true,
    columns=fixed,        % Mantiene alineadas las columnas de Pandas
    keepspaces=true,
    showstringspaces=false,
    literate=
      {á}{{\'a}}1 {é}{{\'e}}1 {í}{{\'i}}1 {ó}{{\'o}}1 {ú}{{\'u}}1
      {Á}{{\'A}}1 {É}{{\'E}}1 {Í}{{\'I}}1 {Ó}{{\'O}}1 {Ú}{{\'U}}1
      {¿}{{?`}}1 {¡}{{!`}}1 {ü}{{\char252}}1 {«}{{\guillemotleft}}1 {»}{{\guillemotright}}1 {·}{{\textperiodcentered}}1 {ñ}{{\~n}}1 {→}{{$\rightarrow$}}1 {Ñ}{{\~N}}1
  },
  title={Salida \thetcbcounter: #2},
  #1
}

%% ---- DATOS DE LA MATERIA ------------------------

\materia{Inteligencia Artificial y su Ética}
\actividad{Tema 6.1. Protección de datos personales \\ Actividad 31 - Datos personales y datos sensibles}
\profesor{M. Lilia Ureña Lugo}


\usepackage{etoolbox}
% ltablex convierte tabularx en longtable, que avanza el contador de tablas una vez de más
\AtEndEnvironment{tabularx}{\addtocounter{table}{-1}}
\newlength{\anchocol}


\begin{document}
\portada
\indicegeneral

\section{Introducción}

Los algoritmos de inteligencia artificial se alimentan de grandes cantidades de datos sobre las personas, y con ellos pueden decidir, por ejemplo, si se aprueba o se niega una tarjeta de crédito sin que un ser humano revise la solicitud \parencite{iae_u6_t61}. Por eso el tema 6.1 pide entender qué se protege y por qué.

Este resumen explica con mis palabras qué entiendo por protección de datos personales y por protección de datos personales sensibles, en qué se distinguen y qué ventajas y desventajas tiene proteger unos y otros. Me apoyo en el material del tema \parencite{iae_u6_t61} y en los textos legales que el tema recomienda leer: la Ley General de Protección de Datos Personales en Posesión de Sujetos Obligados \parencite{lgpdppso2025}, el Reglamento de la Ley Federal de Protección de Datos Personales en Posesión de los Particulares \parencite{reglamento2011} y el documento de la Comisión Nacional de Bioética sobre datos de salud \parencite{conbioetica2014}.

\subsection*{Nota sobre las versiones de las leyes}

El tema lista la ley de sujetos obligados de 2017 y el reglamento de 2011. Revisé los textos oficiales publicados por la Cámara de Diputados y encontré que la ley de 2017 y la ley federal de 2010 fueron abrogadas por el decreto del 20 de marzo de 2025, que expidió una nueva Ley General y una nueva Ley Federal \parencite{lgpdppso2025}. Por eso cito el texto vigente de la Ley General \parencite{lgpdppso2025} y, para el sector privado, la nueva Ley Federal \parencite{lfpdppp2025}, que el tema no lista pero que el reglamento necesita para entenderse. El documento de CONBIOÉTICA es de 2014 y habla de la ley de 2010, así que lo uso por sus ideas sobre salud y no como texto legal vigente.

\section{Qué entiendo por protección de datos personales}

Un \textbf{dato personal} es cualquier información sobre una persona identificada o identificable, es decir, aquella cuya identidad se puede determinar directa o indirectamente \parencite[art.~3, fr.~IX]{lgpdppso2025}. Un nombre, un teléfono o un correo son ejemplos; también lo es un dato que, junto con otros, permite saber quién es la persona.

La \textbf{protección de datos personales} es el conjunto de reglas que obligan a quien trata esos datos (el responsable) a hacerlo bajo ocho principios: licitud, consentimiento, información, calidad, finalidad, lealtad, proporcionalidad y responsabilidad \parencite{iae_u6_t61,lgpdppso2025}. En pocas palabras: se pueden usar mis datos solo para lo que me avisaron, con mi autorización cuando corresponde, usando únicamente los datos necesarios, y el responsable responde si los usa mal. Además, la persona titular tiene los derechos ARCO (acceso, rectificación, cancelación y oposición) \parencite[art.~3, fr.~XI]{lgpdppso2025}, y el responsable debe mantener medidas de seguridad administrativas, técnicas y físicas \parencite[art.~25]{lgpdppso2025}.

\section{Qué entiendo por protección de datos sensibles}

Los \textbf{datos personales sensibles} son los que se refieren a la esfera más íntima de la persona o cuyo mal uso puede causar discriminación o un riesgo grave. La ley da como ejemplos el origen racial o étnico, el estado de salud presente o futuro, la información genética, las creencias religiosas, filosóficas y morales, las opiniones políticas y la preferencia sexual, y aclara que la lista no es limitativa \parencite[art.~3, fr.~X]{lgpdppso2025}. El tema añade los datos biométricos y la afiliación sindical \parencite{iae_u6_t61}.

La \textbf{protección de datos sensibles} usa las mismas reglas que la de datos personales, pero más estrictas. Entiendo la diferencia así: en los datos personales el riesgo es sobre todo la privacidad; en los sensibles, el riesgo es que me traten distinto por lo que soy o por mi salud. Los textos lo reflejan en lo siguiente:

\begin{itemize}
    \item \textbf{Consentimiento.} En el sector privado, el consentimiento tácito es válido por regla general, pero para los datos sensibles se necesita consentimiento expreso y por escrito, con firma autógrafa, electrónica o un mecanismo de autenticación \parencite[arts.~7 y 8]{lfpdppp2025}. En el sector público, por regla general no pueden tratarse datos sensibles sin consentimiento expreso, salvo los casos que prevé la ley \parencite[art.~7]{lgpdppso2025}.
    \item \textbf{Bases de datos.} No pueden crearse bases con datos sensibles si no se justifican para finalidades legítimas, concretas y acordes con las actividades del responsable \parencite[art.~8]{lfpdppp2025}; el reglamento añade que solo se crean por mandato legal o para esas finalidades \parencite[art.~56]{reglamento2011}.
    \item \textbf{Sanciones.} Las multas por infracciones con datos sensibles pueden llegar al doble, y las penas de prisión por tratamiento indebido se duplican \parencite[arts.~59 y 64]{lfpdppp2025}.
\end{itemize}

Los \textbf{datos de salud} son el caso que más trabaja el tema. CONBIOÉTICA parte de que son sensibles y advierte que su divulgación indebida puede causar un daño moral y, sobre todo, discriminación o exclusión, lo que va contra la igualdad entre las personas \parencite{conbioetica2014}. En salud, el principio central es la confidencialidad \parencite{conbioetica2014,iae_u6_t61}.

\section{Ventajas y desventajas de la protección}

\begin{table}[H]
\centering
\scriptsize
\renewcommand{\arraystretch}{1.25}
\begin{tabularx}{\textwidth}{@{}p{2.6cm}XX@{}}
\toprule
 & \textbf{Ventajas} & \textbf{Desventajas o límites} \\
\midrule
\textbf{Datos personales} &
La persona controla sus datos con los derechos ARCO, que son gratuitos \parencite[art.~44]{lgpdppso2025}. Debe recibir un aviso de privacidad \parencite[art.~14]{lfpdppp2025}. Hay medidas de seguridad obligatorias y aviso inmediato cuando una vulneración afecta de forma significativa sus derechos \parencite[arts.~18 y 19]{lfpdppp2025}. Puede oponerse a decisiones tomadas solo por un sistema automatizado con efectos significativos \parencite[art.~26]{lfpdppp2025}. &
En el sector privado el consentimiento tácito basta por regla general \parencite[art.~7]{lfpdppp2025}, y hay varias excepciones al consentimiento, como los datos en fuentes de acceso público o los que se disocian \parencite[art.~9]{lfpdppp2025}. Cumplir implica costos para el responsable (aviso de privacidad, medidas y documento de seguridad); esto es una inferencia mía a partir de las obligaciones, que las fuentes no cuantifican. \\
\addlinespace
\textbf{Datos sensibles} &
Consentimiento expreso y por escrito. Mayor sanción si se usan mal. Se reduce el riesgo de discriminación y exclusión por salud, origen étnico, creencias u orientación sexual \parencite{conbioetica2014}. En salud se refuerza la confidencialidad del expediente clínico y de las muestras biológicas \parencite{iae_u6_t61}. &
Puede chocar con la atención médica y la investigación; CONBIOÉTICA pide proteger la intimidad sin entorpecer la labor sanitaria. En salud, el consentimiento puede omitirse cuando es necesario para la atención médica y la persona no puede darlo, y el propio documento recoge la crítica de que esas excepciones se pueden interpretar tan ampliamente que cubran casi cualquier actuación médica \parencite{conbioetica2014}. El tema señala que muchas disposiciones sanitarias solo enuncian la confidencialidad sin definir su alcance \parencite{iae_u6_t61}. \\
\bottomrule
\end{tabularx}
\caption{Ventajas y desventajas de proteger los datos personales y los sensibles.}\label{tab:ventajas}
\end{table}

\section{Conclusión}

Entiendo la protección de datos personales como el control que cada persona conserva sobre la información que la identifica, y la de datos sensibles como una versión reforzada para los datos cuyo mal uso puede discriminarla. Las ventajas son el control del titular, la seguridad obligatoria y las sanciones. Las desventajas están en los límites de la norma (consentimiento tácito, excepciones) y en el equilibrio con la atención médica y la investigación. Para quien diseña sistemas de IA esto importa porque los datos de entrenamiento suelen ser datos personales, y a veces sensibles, y el tema plantea que hay que fijar límites al uso de datos y algoritmos \parencite{iae_u6_t61}.

\subsection*{Limitaciones}

Los libros que el tema cita (Kearns y Roth; Coeckelbergh; la CIDH) no los leí, y solo uso lo que el tema dice de ellos. El documento de CONBIOÉTICA es un escaneo de 2014; lo leí con reconocimiento de texto y citó leyes que ya no están vigentes. No revisé jurisprudencia ni casos de aplicación de las leyes.

\referencias

\end{document}
